\documentclass[type=paper,theme=light,language=de,citestyle=alphabetic]{semio}
\title{Probe2}
\author{Probe}
\date{2026}

\makeatletter
\newcommand{\ProbeDump}[1]{%
  \typeout{[DEBUG] #1 parindent=\the\parindent}%
  \typeout{[DEBUG] #1 leftskip=\the\leftskip}%
  \typeout{[DEBUG] #1 hsize=\the\hsize}%
  \typeout{[DEBUG] #1 cellpad=\the\semio@table@long@cellpad}%
  \typeout{[DEBUG] #1 innerw=\the\semio@table@long@inner@w}%
  \typeout{[DEBUG] #1 hpad=\semio@table@hpad}%
}
\makeatother

\begin{document}

\section{Probe}

\begin{Table}[title=Current Cellfill]
\begin{SemioTable}{0.20,0.10,0.18,0.27,0.25}
  \SemioTableHeaderRow{Risiko & Status & Verantwortung & Maßnahme & Restrisiko}
  \SemioTableRow{\ProbeDump{col1}Formale Fristen und interne Meilensteinplanung verwenden unterschiedliche Ausgangspunkte & adressiert & \ProbeDump{col3}NGS, Verbundkoordination & Formaler und operativer Projektbeginn sind mit ihren jeweiligen Fristen im Bericht getrennt dargestellt & Formale Bewertung kann von der internen Zeitrechnung abweichen}
\end{SemioTable}
\end{Table}

\makeatletter
\renewcommand{\semio@table@long@cellfill}{%
  \cellcolor{semio-chrome-canvas}%
  \raggedright\hyphenpenalty=50\exhyphenpenalty=50\emergencystretch=1.2em\arraybackslash\hspace{0pt}\semio@cell@struttop
}
\renewcommand{\semio@table@long@cellfill@first}{%
  \cellcolor{semio-chrome-canvas}%
  \raggedright\hyphenpenalty=50\exhyphenpenalty=50\emergencystretch=1.2em\arraybackslash\hspace{0pt}\semio@cell@struttop
}
\makeatother

\begin{Table}[title=Glue First No Leftskip]
\begin{SemioTable}{0.20,0.10,0.18,0.27,0.25}
  \SemioTableHeaderRow{Risiko & Status & Verantwortung & Maßnahme & Restrisiko}
  \SemioTableRow{Formale Fristen und interne Meilensteinplanung verwenden unterschiedliche Ausgangspunkte & adressiert & NGS, Verbundkoordination & Formaler und operativer Projektbeginn sind mit ihren jeweiligen Fristen im Bericht getrennt dargestellt & Formale Bewertung kann von der internen Zeitrechnung abweichen}
\end{SemioTable}
\end{Table}

\makeatletter
\renewcommand{\semio@table@long@cellfill}{%
  \cellcolor{semio-chrome-canvas}%
  \raggedright\hyphenpenalty=50\exhyphenpenalty=50\emergencystretch=1.2em\arraybackslash\semio@cell@struttop\hspace{0pt}\ignorespaces
}
\renewcommand{\semio@table@long@cellfill@first}{%
  \cellcolor{semio-chrome-canvas}%
  \raggedright\hyphenpenalty=50\exhyphenpenalty=50\emergencystretch=1.2em\arraybackslash\semio@cell@struttop\hspace{0pt}\ignorespaces
}
\makeatother

\begin{Table}[title=Strut First Ignorespaces]
\begin{SemioTable}{0.20,0.10,0.18,0.27,0.25}
  \SemioTableHeaderRow{Risiko & Status & Verantwortung & Maßnahme & Restrisiko}
  \SemioTableRow{Formale Fristen und interne Meilensteinplanung verwenden unterschiedliche Ausgangspunkte & adressiert & NGS, Verbundkoordination & Formaler und operativer Projektbeginn sind mit ihren jeweiligen Fristen im Bericht getrennt dargestellt & Formale Bewertung kann von der internen Zeitrechnung abweichen}
\end{SemioTable}
\end{Table}

\end{document}
